% Рубежный контроль №1 «Моделирование» — Вариант 1
% Источники: mathmod.pdf, Modelirovanie_1_-3.pdf, 23_03_2020_..._Моделирование_2.doc
\documentclass[12pt,a4paper]{article}
\usepackage[T2A]{fontenc}
\usepackage[utf8]{inputenc}
\usepackage[russian]{babel}
\usepackage{amsmath,amssymb}
\usepackage{geometry}
\geometry{margin=2.2cm}
\usepackage{enumitem}
\setlist{nosep,leftmargin=*}

\title{Рубежный контроль №1 «Моделирование»\\Вариант 1 — ответы}
\author{}
\date{}

\begin{document}
\maketitle

\begin{enumerate}[label=\textbf{\arabic*.}]

\item \textbf{Аналитическое моделирование (определение).}

Аналитическая модель — формализация явления/процесса реального мира, учитывающая физико-химические воздействия, протекающие в системе, и задающая связи между величинами в аналитическом (формульном) виде. Аналитическое моделирование — построение и исследование таких моделей (в т.\,ч. получение решений прямых/обратных задач и анализ адекватности).

\item \textbf{Имитационная модель (определение).}

Имитационная модель — численный метод вычислительного эксперимента, позволяющий имитировать поведение реальной системы на ЦВМ в заданный или формируемый период времени при заданных входных данных (и, при необходимости, управляющих воздействиях). Для многокомпонентных систем имитация опирается на дискретное модельное время и механизмы синхронизации компонент.

\item \textbf{Полунатурный эксперимент (определение).}

Полунатурный эксперимент — компромисс между натурным и вычислительным: объектом исследования является физическая модель реального объекта; часть характеристик измеряется экспериментально, а часть рассчитывается программно/по математической модели (возможны итерации уточнения параметров).

\item \textbf{Определение модельного времени.}

Модельное время — дискретное время, измеряемое в единицах реального времени, используемое для внутренней и внешней синхронизации вычислительных процессов (компонент имитационной модели). Один квант модельного времени соответствует переходу системы из одного состояния в другое.

\item \textbf{Определение сложной динамической системы, меняющей поведение во времени.}

Сложная динамическая система, меняющая поведение во времени — система, для которой в разные интервалы времени реализуются различные режимы (законы) функционирования: меняются параметры, ограничения, структура связей между компонентами или алгоритмы преобразования входов в выходы. Для описания часто требуется декомпозиция по режимам/уровням и исследование переходных функций на границах режимов; возможны разрывы и гибридность поведения.

\item \textbf{Способы организации квазипараллелизма (перечислить).}

Схемы организации квазипараллелизма (псевдопараллелизма) при дискретном модельном времени:
\begin{enumerate}[label=\arabic*)]
\item активностями;
\item событиями;
\item транзактами;
\item процессами;
\item агрегатами.
\end{enumerate}

\item \textbf{Построение модифицированного метода моментов (линейная функциональная зависимость).}

\textbf{Постановка.} Даны пары наблюдений $(x_i,y_i)$, $i=1,\dots,n$, и предполагается линейная связь (линейная регрессия)
\[
\mathbb{E}[Y\mid X=x]=\theta_0+\theta_1 x.
\]
Оценки параметров связи можно получить как \textbf{модификацию метода моментов}: приравнять теоретические моменты (первый момент $Y$ и совместный момент/ковариацию $(X,Y)$) их выборочным аналогам:
\[
\mathbb{E}[Y]=\theta_0+\theta_1\mathbb{E}[X], \qquad \mathrm{cov}(X,Y)=\theta_1\,\mathrm{var}(X).
\]
Заменяя ожидания/дисперсии/ковариацию выборочными величинами
\[
\bar x=\frac1n\sum_{i=1}^n x_i,\quad \bar y=\frac1n\sum_{i=1}^n y_i,\quad
\widehat S_{xx}=\frac1n\sum_{i=1}^n (x_i-\bar x)^2,\quad
\widehat S_{xy}=\frac1n\sum_{i=1}^n (x_i-\bar x)(y_i-\bar y),
\]
получаем
\[
\hat\theta_1=\frac{\widehat S_{xy}}{\widehat S_{xx}}, \qquad
\hat\theta_0=\bar y-\hat\theta_1\bar x,
\]
при $\widehat S_{xx}>0$. (По форме совпадает с оценками МНК для простой линейной регрессии.)

\item \textbf{Функциональная зависимость (определение).}

Функциональная зависимость — соответствие, при котором каждой величине $x\in X$ ставится в соответствие единственное значение $y\in Y$ (в отличие от стохастической зависимости, где одному $x$ могут соответствовать различные $y$).

\item \textbf{9–10. Имеются результаты ЕГЭ 10 абитуриентов. Оценить силу связи между показателями посредством рангового коэффициента Спирмена.}

Для двух показателей $X$ и $Y$ ранговый коэффициент Спирмена вычисляют как корреляцию Пирсона между рангами:
\[
\rho_s=\mathrm{corr}(\mathrm{rank}(X),\mathrm{rank}(Y)).
\]
При наличии совпадающих значений используются усреднённые ранги.

\textbf{Связь между РЯ и М.}

\textbf{Шаг 1. Ранги (по возрастанию баллов; при совпадениях — средний ранг).}

Для РЯ (82, 78, 85, 96, 98, 56, 84, 88, 72, 66) получаем ранги:
\[
r_{\text{РЯ}}=(5,\ 4,\ 7,\ 9,\ 10,\ 1,\ 6,\ 8,\ 3,\ 2).
\]

Для М (76, 82, 90, 98, 96, 76, 78, 98, 80, 82) ранги с учётом повторов:
\[
r_{\text{М}}=(1.5,\ 5.5,\ 7,\ 9.5,\ 8,\ 1.5,\ 3,\ 9.5,\ 4,\ 5.5).
\]

\textbf{Шаг 2. Корреляция Пирсона для рангов.}
Обозначим $\bar r_X=\frac1n\sum r_{X,i}$, $\bar r_Y=\frac1n\sum r_{Y,i}$ (здесь $n=10$, поэтому $\bar r_X=\bar r_Y=5.5$).
Тогда
\[
\rho_s=\frac{\sum_{i=1}^n (r_{X,i}-\bar r_X)(r_{Y,i}-\bar r_Y)}
{\sqrt{\sum_{i=1}^n (r_{X,i}-\bar r_X)^2}\ \sqrt{\sum_{i=1}^n (r_{Y,i}-\bar r_Y)^2}}.
\]
Вычисления дают:
\[
\sum (r_{\text{РЯ}}-5.5)(r_{\text{М}}-5.5)=60.0,\quad
\sum (r_{\text{РЯ}}-5.5)^2=82.5,\quad
\sum (r_{\text{М}}-5.5)^2=81.0,
\]
откуда
\[
\rho_s(\text{РЯ},\text{М})=\frac{60.0}{\sqrt{82.5\cdot 81.0}}\approx 0.734.
\]

\textbf{Ответ:} $\rho_s(\text{РЯ},\text{М})\approx 0.734$.
\[
\rho_s(\text{РЯ},\text{М})\approx 0.734.
\]

\textbf{Связь между Ф и И.}

\textbf{Шаг 1. Ранги.}

Для Ф (85, 80, 88, 92, 100, 74, 72, 90, 65, 69):
\[
r_{\text{Ф}}=(7,\ 5,\ 8,\ 9,\ 10,\ 4,\ 3,\ 6,\ 1,\ 2).
\]

Для И (68, 82, 94, 100, 84, 64, 70, 90, 74, 78):
\[
r_{\text{И}}=(2,\ 7,\ 9,\ 10,\ 8,\ 1,\ 3,\ 6,\ 4,\ 5).
\]

\textbf{Шаг 2. Корреляция рангов.}
\[
\sum (r_{\text{Ф}}-5.5)(r_{\text{И}}-5.5)=53.5,\quad
\sum (r_{\text{Ф}}-5.5)^2=82.5,\quad
\sum (r_{\text{И}}-5.5)^2=82.5,
\]
поэтому
\[
\rho_s(\text{Ф},\text{И})=\frac{53.5}{82.5}\approx 0.648.
\]

\textbf{Ответ:} $\rho_s(\text{Ф},\text{И})\approx 0.648$.
\[
\rho_s(\text{Ф},\text{И})\approx 0.648.
\]

\textbf{Связь между остальными парами предметов (аналогично, по рангам).}

\begin{itemize}
\item \textbf{РЯ и Ф.} Ранги $r_{\text{РЯ}}=(5,4,7,9,10,1,6,8,3,2)$, $r_{\text{Ф}}=(6,5,7,9,10,4,3,8,1,2)$.
\[
\sum (r_{\text{РЯ}}-5.5)(r_{\text{Ф}}-5.5)=70.5,\quad
\sum (r_{\text{РЯ}}-5.5)^2=82.5,\quad
\sum (r_{\text{Ф}}-5.5)^2=82.5,
\]
\[
\rho_s(\text{РЯ},\text{Ф})=\frac{70.5}{82.5}\approx 0.855.
\]
\item \textbf{РЯ и И.} Ранги $r_{\text{РЯ}}=(5,4,7,9,10,1,6,8,3,2)$, $r_{\text{И}}=(2,6,9,10,7,1,3,8,4,5)$.
\[
\sum (r_{\text{РЯ}}-5.5)(r_{\text{И}}-5.5)=59.5,\quad
\sum (r_{\text{РЯ}}-5.5)^2=82.5,\quad
\sum (r_{\text{И}}-5.5)^2=82.5,
\]
\[
\rho_s(\text{РЯ},\text{И})=\frac{59.5}{82.5}\approx 0.721.
\]
\item \textbf{М и Ф.} Ранги $r_{\text{М}}=(1.5,5.5,7,9.5,8,1.5,3,9.5,4,5.5)$, $r_{\text{Ф}}=(6,5,7,9,10,4,3,8,1,2)$.
\[
\sum (r_{\text{М}}-5.5)(r_{\text{Ф}}-5.5)=54.5,\quad
\sum (r_{\text{М}}-5.5)^2=81.0,\quad
\sum (r_{\text{Ф}}-5.5)^2=82.5,
\]
\[
\rho_s(\text{М},\text{Ф})=\frac{54.5}{\sqrt{81.0\cdot 82.5}}\approx 0.667.
\]
\item \textbf{М и И.} Ранги $r_{\text{М}}=(1.5,5.5,7,9.5,8,1.5,3,9.5,4,5.5)$, $r_{\text{И}}=(2,6,9,10,7,1,3,8,4,5)$.
\[
\sum (r_{\text{М}}-5.5)(r_{\text{И}}-5.5)=77.5,\quad
\sum (r_{\text{М}}-5.5)^2=81.0,\quad
\sum (r_{\text{И}}-5.5)^2=82.5,
\]
\[
\rho_s(\text{М},\text{И})=\frac{77.5}{\sqrt{81.0\cdot 82.5}}\approx 0.948.
\]
\end{itemize}

\textbf{Итог по всем 6 парам.}
\[
\rho_s(\text{РЯ},\text{М})\approx 0.734,\quad
\rho_s(\text{РЯ},\text{Ф})\approx 0.855,\quad
\rho_s(\text{РЯ},\text{И})\approx 0.721,
\]
\[
\rho_s(\text{М},\text{Ф})\approx 0.667,\quad
\rho_s(\text{М},\text{И})\approx 0.948,\quad
\rho_s(\text{Ф},\text{И})\approx 0.648.
\]
\textbf{Вывод:} во всех парах связь \emph{положительная}. Самая сильная — между \textbf{М} и \textbf{И} ($\approx 0.948$), затем \textbf{РЯ} и \textbf{Ф} ($\approx 0.855$). Самая слабая из шести — между \textbf{Ф} и \textbf{И} ($\approx 0.648$) и между \textbf{М} и \textbf{Ф} ($\approx 0.667$).

\end{enumerate}
\end{document}

